\documentclass[11pt,a4paper]{report}

% ---------- Encodage et langue ----------
\usepackage[utf8]{inputenc}
\usepackage[T1]{fontenc}
\usepackage[french]{babel}
\usepackage{lmodern}

% ---------- Mise en page ----------
\usepackage[margin=2.5cm]{geometry}
\usepackage{setspace}
\setstretch{1.15}
\usepackage{titlesec}
\usepackage{fancyhdr}
\pagestyle{fancy}
\fancyhf{}
\fancyfoot[C]{\thepage}
\renewcommand{\headrulewidth}{0pt}

% ---------- Couleurs ----------
\usepackage{xcolor}
\definecolor{navy}{HTML}{12213D}
\definecolor{bluetitle}{HTML}{1D3A6E}
\definecolor{grey}{HTML}{4B5563}
\definecolor{codebg}{HTML}{0F172A}
\definecolor{codefg}{HTML}{E2E8F0}
\definecolor{calloutbg}{HTML}{EFF6FF}
\definecolor{warnbg}{HTML}{FEF2F2}
\definecolor{warnred}{HTML}{DC2626}

% ---------- Titres ----------
\titleformat{\chapter}[hang]{\normalfont\huge\bfseries\color{navy}}{\thechapter.}{12pt}{}
\titleformat{\section}{\normalfont\Large\bfseries\color{navy}}{\thesection}{8pt}{}
\titleformat{\subsection}{\normalfont\large\bfseries\color{bluetitle}}{\thesubsection}{8pt}{}

% ---------- Liens et sommaire ----------
\usepackage[hidelinks]{hyperref}
\usepackage{tocloft}

% ---------- Tableaux ----------
\usepackage{array}
\usepackage{longtable}
\usepackage{booktabs}
\usepackage{colortbl}

% ---------- Blocs de code ----------
\usepackage{listings}
\lstdefinestyle{codebox}{
  backgroundcolor=\color{codebg},
  basicstyle=\ttfamily\footnotesize\color{codefg},
  breaklines=true,
  frame=none,
  xleftmargin=8pt,
  xrightmargin=8pt,
  framexleftmargin=8pt,
  framexrightmargin=8pt,
  framesep=8pt,
}
\lstset{style=codebox}

% ---------- Encadrés (callouts) ----------
\usepackage{tcolorbox}
\tcbuselibrary{skins,breakable}

\newtcolorbox{infobox}{
  colback=calloutbg, colframe=calloutbg, boxrule=0pt,
  leftrule=3pt, colframe=bluetitle,
  arc=1pt, boxsep=2pt, left=8pt, right=8pt, top=6pt, bottom=6pt,
  breakable
}
\newtcolorbox{warnbox}{
  colback=warnbg, colframe=warnbg, boxrule=0pt,
  leftrule=3pt, colframe=warnred,
  arc=1pt, boxsep=2pt, left=8pt, right=8pt, top=6pt, bottom=6pt,
  breakable
}

\usepackage{enumitem}
\setlist{itemsep=2pt, topsep=4pt}

\begin{document}

% ================= SOMMAIRE =================
\renewcommand{\contentsname}{Sommaire}
\tableofcontents
\clearpage

% ================= 1. INTRODUCTION =================
\chapter{Introduction générale}

Dans le cadre de ma formation, j'ai eu l'opportunité d'effectuer un stage portant sur la conception, le
développement et la consolidation d'une plateforme web nommée \textbf{SARAI} (\emph{Stocktaking of Arab
Regional AI Initiatives}), destinée à recenser et cartographier les initiatives d'intelligence artificielle
dans la région arabe pour le compte de l'\textbf{AICTO} (Arab ICT Organization), organisation régionale
œuvrant pour le développement des technologies de l'information et de la communication dans le monde arabe.

Ce stage m'a permis de travailler sur une application full-stack complète, combinant une API backend
construite avec le framework \textbf{FastAPI} (Python), une base de données relationnelle
\textbf{PostgreSQL} enrichie de capacités de recherche vectorielle (\texttt{pgvector}), et une interface
utilisateur moderne développée avec \textbf{React} et \textbf{Vite}. Le projet couvre un périmètre
fonctionnel riche~: gestion des utilisateurs et des rôles, soumission et modération de projets IA,
cartographie interactive par pays, tableau de bord statistique, moteur de recherche hybride (mots-clés,
sémantique, plein texte), chatbot conversationnel, notifications, génération de rapports PDF, et prise en
charge de trois langues (français, anglais, arabe) avec gestion du sens d'écriture bidirectionnel.

Ce rapport présente dans un premier temps le contexte et les objectifs du projet, puis détaille les choix
technologiques et l'architecture retenue, avant de décrire les principales réalisations effectuées durant
le stage. Une attention particulière est portée aux travaux d'audit, de sécurisation et de nettoyage du
code, qui ont constitué une part importante de la mission. Le rapport se conclut par un bilan des
compétences acquises et des perspectives d'évolution du projet.

% ================= 2. CONTEXTE =================
\chapter{Contexte du projet}

\section{Présentation de l'organisme d'accueil}

L'AICTO (Arab ICT Organization) est une organisation régionale dont la mission est de promouvoir la
coopération entre les pays arabes dans le domaine des technologies de l'information et de la communication.
Dans ce cadre, l'organisation a souhaité disposer d'un outil numérique permettant de dresser un état des
lieux (« stocktaking ») des initiatives d'intelligence artificielle menées dans les 22 pays de la région
arabe, afin de favoriser le partage de connaissances, la mise en réseau des acteurs (startups, universités,
laboratoires de recherche, agences gouvernementales, ONG, entreprises) et l'alignement de ces initiatives
avec les Objectifs de Développement Durable (ODD) des Nations Unies.

\section{Problématique et objectifs}

Avant SARAI, les initiatives IA de la région n'étaient recensées dans aucun système centralisé~:
l'information était dispersée entre rapports, sites institutionnels et bases de données locales non
interopérables. Le projet répond donc aux objectifs suivants~:

\begin{itemize}
  \item Centraliser les données relatives aux projets IA, aux parties prenantes (organisations), aux
    secteurs d'activité, aux technologies IA employées et aux ODD associés~;
  \item Offrir une cartographie interactive permettant de visualiser la répartition géographique des
    initiatives~;
  \item Fournir un tableau de bord analytique pour le suivi des tendances (par pays, secteur, technologie,
    statut)~;
  \item Garantir la qualité des données via un workflow de modération (soumission puis validation par un
    administrateur)~;
  \item Rendre la plateforme accessible en français, anglais et arabe, compte tenu du public régional
    visé~;
  \item Assister les utilisateurs via un chatbot capable de répondre à des requêtes en langage naturel.
\end{itemize}

\section{Cahier des charges fonctionnel (synthèse)}

\begin{longtable}{@{}p{4cm}p{10.5cm}@{}}
\toprule
\textbf{Module} & \textbf{Fonctionnalités attendues} \\
\midrule
\endhead
Authentification & Inscription, vérification d'email, connexion, mot de passe oublié, gestion de rôles (organisation, modérateur, administrateur) \\
Gestion de projets & Soumission, édition, association à un pays/secteur/technologies/ODD, statut de workflow, historique d'approbation \\
Annuaire des parties prenantes & Fiche organisation, type, pays, secteur \\
Cartographie & Carte interactive par pays avec indicateurs socio-économiques et statuts de projets \\
Statistiques & Tableau de bord agrégé (répartition par pays, secteur, technologie, statut) \\
Recherche & Recherche par mots-clés/filtres, recherche sémantique, option plein texte (Elasticsearch) \\
Chatbot & Assistant conversationnel multilingue répondant sur la base des données de la plateforme \\
Modération & File d'attente admin, approbation/rejet individuel ou en masse, notifications \\
Ressources & Bibliothèque de rapports, jeux de données, livres blancs \\
Export & Génération de fiches projet au format PDF \\
Internationalisation & FR / EN / AR avec bascule RTL/LTR automatique \\
\bottomrule
\end{longtable}

% ================= 3. ETUDE TECHNO =================
\chapter{Étude et choix technologiques}

\section{Architecture générale}

Le projet adopte une architecture \textbf{monolithique modulaire} plutôt qu'une architecture en
microservices. Ce choix, documenté explicitement dans la documentation d'architecture du projet, repose sur
quatre arguments~:

\begin{enumerate}
  \item le domaine métier est fortement relationnel (jointures fréquentes entre projets, parties prenantes,
    pays et ODD), ce qui rend un découpage en microservices coûteux sans bénéfice réel~;
  \item la charge de la plateforme est modérée et ne justifie pas un scaling indépendant des composants~;
  \item l'équipe de développement est restreinte, rendant la coordination de plusieurs services distants
    contre-productive~;
  \item les seules « coutures » naturelles du système — la recherche (Elasticsearch), les embeddings
    sémantiques (pgvector) et le chatbot IA (Ollama / API compatible OpenAI) — sont déjà isolées dans des
    modules de services dédiés et pourront être extraites plus tard sans toucher au cœur métier.
\end{enumerate}

\begin{infobox}
\textbf{Point d'architecture retenu~:} ce raisonnement correspond à la démarche d'un \emph{Architecture
Decision Record} (ADR) — une décision technique justifiée et documentée plutôt qu'un choix par défaut. Ce
type de réflexion a constitué un apprentissage important durant le stage.
\end{infobox}

\section{Technologies côté backend}

\begin{longtable}{@{}p{3cm}p{4.5cm}p{7cm}@{}}
\toprule
\textbf{Composant} & \textbf{Technologie} & \textbf{Rôle} \\
\midrule
\endhead
Framework API & FastAPI 0.109 & Exposition des endpoints REST, validation automatique, documentation Swagger générée \\
Serveur ASGI & Uvicorn & Exécution asynchrone de l'application \\
ORM & SQLAlchemy 2.0 & Mapping objet-relationnel \\
Validation & Pydantic 2.5 & Schémas de validation des requêtes/réponses (DTO) \\
Base de données & PostgreSQL 16 (\texttt{pgvector/pgvector:pg16}) & Persistance relationnelle + recherche vectorielle, repli automatique sur SQLite en développement \\
Authentification & JWT (\texttt{python-jose}), \texttt{passlib} (pbkdf2\_sha256) & Jetons d'accès, hachage sécurisé des mots de passe \\
Recherche & Elasticsearch (optionnel), \texttt{sentence-transformers} + pgvector & Recherche plein texte et recherche sémantique par similarité d'embeddings \\
Chatbot & Ollama (llama3.2, self-hosted) ou API compatible OpenAI & Assistant conversationnel \\
Génération PDF & \texttt{xhtml2pdf} + Jinja2 & Export des fiches projet en PDF, avec résumé généré par IA \\
Emails & SMTP (service dédié) & Vérification d'email, réinitialisation de mot de passe, notifications \\
\bottomrule
\end{longtable}

\section{Technologies côté frontend}

\begin{longtable}{@{}p{3.5cm}p{4cm}p{7cm}@{}}
\toprule
\textbf{Composant} & \textbf{Technologie} & \textbf{Rôle} \\
\midrule
\endhead
Bibliothèque UI & React 18 & Construction de l'interface en composants \\
Outil de build & Vite 5 & Bundling rapide, rechargement à chaud en développement \\
Routage & \texttt{react-router-dom} v6 & Navigation SPA \\
Client HTTP & \texttt{axios} & Appels à l'API backend \\
Internationalisation & \texttt{i18next} / \texttt{react-i18next} & Traduction FR / EN / AR, détection de langue \\
Visualisation & \texttt{recharts} & Graphiques du tableau de bord analytique \\
Notifications UI & \texttt{react-toastify} & Messages toast avec support RTL \\
Export PDF client & \texttt{jspdf}, \texttt{html2canvas} & Génération de documents côté navigateur \\
\bottomrule
\end{longtable}

\section{Justification des choix}

FastAPI a été retenu pour sa performance (asynchrone), sa validation de schémas intégrée via Pydantic, et sa
documentation interactive générée automatiquement (Swagger UI), particulièrement utile pour la collaboration
frontend/backend. PostgreSQL avec l'extension pgvector permet de combiner données relationnelles classiques
et recherche par similarité vectorielle dans un seul système, évitant l'introduction d'une base de données
spécialisée supplémentaire. Côté frontend, React et Vite offrent un environnement de développement rapide et
un écosystème mature pour couvrir les besoins d'internationalisation, de theming et de visualisation de
données.

% ================= 4. CONCEPTION =================
\chapter{Conception}

\section{Architecture logicielle en couches (backend)}

Le backend est organisé selon une architecture en couches, séparant clairement les responsabilités~:

\begin{lstlisting}
backend/app/
|-- routers/       -> couche HTTP : endpoints REST (pas de logique metier)
|-- services/      -> logique metier + integrations externes (ES, IA, mail, PDF)
|-- models/        -> entites SQLAlchemy (persistance en base de donnees)
|-- schemas/       -> DTO Pydantic (contrat de donnees a la frontiere HTTP)
`-- dependencies.py -> authentification / autorisation transversale (JWT, RBAC)
\end{lstlisting}
\begin{center}\small\color{grey} Figure~1 --- Organisation en couches du backend FastAPI \end{center}

\section{Modèle de données}

Le schéma de base de données a fait l'objet d'une refonte vers une forme normalisée (troisième forme
normale) comportant 18 tables. Les entités principales sont résumées ci-dessous~:

\begin{longtable}{@{}p{3.5cm}p{10.5cm}@{}}
\toprule
\textbf{Entité} & \textbf{Description} \\
\midrule
\endhead
\texttt{User} & Comptes utilisateurs~: email, mot de passe haché, rôle (organisation / modérateur / administrateur), statuts d'activation, de vérification d'email et d'approbation admin \\
\texttt{Project} & Entité centrale~: projet IA, statut de workflow (en attente, approuvé, rejeté, idée, prototype, production, en cours, terminé), colonne d'embedding vectoriel pour la recherche sémantique \\
\texttt{Stakeholder} & Organisations porteuses de projets (type, pays, secteur) \\
\texttt{Country} & 22 pays arabes avec coordonnées géographiques \\
\texttt{Sector} & Taxonomie des secteurs d'activité \\
\texttt{SDG} & 17 Objectifs de Développement Durable de l'ONU \\
\texttt{AITechnology} & Taxonomie des technologies d'IA (NLP, vision, robotique, LLM, etc.) \\
\texttt{Comment}, \texttt{Approval} & Commentaires modérés et historique d'audit du workflow d'approbation \\
\texttt{Notification} & Notifications utilisateur (approbation, rejet, commentaire, mention, système) \\
\texttt{Resource} & Bibliothèque de ressources documentaires \\
\texttt{VerificationToken}, \texttt{PasswordResetToken} & Jetons pour la vérification d'email et la réinitialisation de mot de passe \\
\bottomrule
\end{longtable}

Des tables de liaison assurent les relations plusieurs-à-plusieurs entre projets et ODD, projets et
technologies, projets et parties prenantes. Le modèle \texttt{Project} intègre également des champs métier
enrichis (couverture géographique, tâches planifiées, impact attendu, budget, durée prévue), ajoutés
progressivement au fil de l'évolution du cahier des charges.

\section{Architecture de déploiement}

\begin{lstlisting}
+----------------------------------------------------------+
|                     Docker Compose                        |
|                                                            |
|  +-----------+   +--------------+   +--------------------+
|  | frontend  |-->|   backend    |-->|         db          |
|  | (Nginx)   |   | (FastAPI /   |   | (PostgreSQL 16 +    |
|  | port 3001 |   |  Uvicorn)    |   |  extension pgvector) |
|  +-----------+   |  port 8000   |   +--------------------+
|                  +------+-------+                          |
|                         | (profils optionnels)              |
|            +------------+-------------+                    |
|    +-------v------+           +--------v---------+          |
|    | elasticsearch |           |      ollama       |          |
|    | (profil       |           | (profil "ai" --   |          |
|    |  "search")    |           |  chatbot local)   |          |
|    +---------------+           +-------------------+          |
+----------------------------------------------------------+
\end{lstlisting}
\begin{center}\small\color{grey} Figure~2 --- Architecture de déploiement (docker-compose.yml) \end{center}

Le service backend utilise une image Python 3.12 slim exécutée sous un utilisateur non privilégié. Le
frontend est construit en deux étapes (build multi-stage)~: compilation Vite dans une image Node 22, puis
service des fichiers statiques via Nginx, avec une configuration dédiée assurant le repli sur
\texttt{index.html} pour le routage côté client (SPA). L'URL de l'API est injectée au moment du build via la
variable d'environnement \texttt{VITE\_API\_URL}, ce qui permet de déployer la même image sur différents
environnements sans recompilation.

% ================= 5. REALISATION =================
\chapter{Réalisation}

\section{Authentification et gestion des rôles}

L'authentification repose sur des jetons JWT. Lors de la connexion, le serveur génère un jeton signé
(HS256) valable 24 heures, contenant l'identité et le rôle de l'utilisateur~:

\begin{lstlisting}
oauth2_scheme = OAuth2PasswordBearer(tokenUrl="/api/auth/login", auto_error=True)

def get_current_user(token: str = Depends(oauth2_scheme), db: Session = Depends(get_db)) -> User:
    payload = jwt.decode(token, SECRET_KEY, algorithms=["HS256"])
    ...

def require_admin(current_user: User = Depends(get_current_user)) -> User:
    if current_user.role != "admin":
        raise HTTPException(status_code=403, ...)
\end{lstlisting}

Trois rôles sont gérés~: organisation, modérateur et administrateur, contrôlés à la fois par une contrainte
SQL au niveau de la base de données et par des dépendances FastAPI appliquées à l'ensemble des routes
sensibles (par exemple, tout le routeur d'administration est protégé par
\texttt{APIRouter(dependencies=[Depends(require\_admin)])}). Le processus de connexion impose également
trois garde-fous métier avant d'autoriser l'accès~: l'email doit être vérifié, le compte doit avoir été
approuvé par un administrateur, et le compte doit être actif.

Côté client, le jeton est conservé dans le \texttt{localStorage} du navigateur puis réinjecté dans l'en-tête
\texttt{Authorization: Bearer <token>} pour chaque appel nécessitant une authentification.

\section{Gestion des projets IA et workflow de modération}

La page de recensement des projets (\texttt{ProjectStocktaking}) constitue le cœur fonctionnel de
l'application~: elle permet à une organisation de soumettre un projet en le rattachant à un pays, un ou
plusieurs secteurs, une ou plusieurs technologies IA et un ou plusieurs Objectifs de Développement Durable.
Chaque soumission entre dans un workflow de modération~: le projet est créé avec le statut « en attente »,
puis un administrateur peut l'approuver, le rejeter (avec motif) ou demander une révision. Chaque changement
de statut est tracé dans une table d'audit dédiée et déclenche une notification (en base et par email) à
destination de l'auteur. Le tableau de bord d'administration permet en outre des actions groupées
(approbation ou rejet en masse) afin de traiter efficacement un grand nombre de soumissions.

\section{Cartographie interactive et statistiques}

La cartographie des initiatives (page \texttt{KnowledgeMap}) affiche une carte des pays arabes avec, pour
chaque pays, des indicateurs socio-économiques (population, PIB, indice de développement humain, devise) et
un code couleur reflétant les statuts des projets qui y sont recensés. Il est à noter que, bien que les
bibliothèques \texttt{leaflet} et \texttt{react-leaflet} figurent parmi les dépendances du projet, la carte a
finalement été implémentée sous la forme d'un composant SVG entièrement personnalisé, avec une projection
géographique codée manuellement — un choix permettant un contrôle visuel fin sur l'affichage des marqueurs
par pays. Le tableau de bord statistique (page \texttt{Analytics}), quant à lui, s'appuie sur des
agrégations calculées côté serveur et restituées sous forme de graphiques via la bibliothèque
\texttt{recharts}.

\section{Recherche et chatbot intelligent}

La plateforme propose une recherche hybride combinant une recherche classique par mots-clés et filtres
(pays, secteur, technologie), une recherche sémantique fondée sur des embeddings générés par
\texttt{sentence-transformers} et comparés via pgvector, ainsi qu'une option de recherche plein texte via
Elasticsearch lorsque ce service est activé. Un chatbot conversationnel est également intégré, alimenté par
un modèle de langage exécuté localement via Ollama (modèle llama3.2) ou, en repli, par une API compatible
OpenAI. Un mécanisme de normalisation d'alias multilingues permet au chatbot de comprendre des requêtes
formulées en français ou en anglais faisant référence aux pays, secteurs ou technologies (par exemple,
« projets IA santé Tunisie »).

\section{Internationalisation et accessibilité}

L'application est disponible en français, anglais et arabe, grâce à \texttt{i18next} et
\texttt{react-i18next}. Au-delà de la simple traduction des textes, le passage à l'arabe déclenche une
bascule complète de la mise en page en mode RTL (droite-à-gauche)~:

\begin{lstlisting}
const dir = i18n.language === 'ar' ? 'rtl' : 'ltr'
<div className="app" dir={dir}>
\end{lstlisting}

Cette prise en charge s'étend jusqu'aux composants tiers, comme les notifications toast dont l'orientation
est également adaptée (\texttt{rtl=\{dir === 'rtl'\}}). Un mode sombre/clair est par ailleurs disponible,
géré par un contexte React persistant dans le \texttt{localStorage} et initialisé automatiquement selon la
préférence système de l'utilisateur (\texttt{prefers-color-scheme}).

\section{Notifications et génération de rapports PDF}

Un système de notifications en base de données informe chaque utilisateur des événements le concernant
(approbation ou rejet de projet, ajout de commentaire, mention, demande de révision), avec un compteur de
messages non lus visible dans la barre de navigation. Par ailleurs, chaque fiche projet peut être exportée
au format PDF~: côté serveur, cette génération s'appuie sur des gabarits Jinja2 associés à la bibliothèque
\texttt{xhtml2pdf}, et peut inclure un résumé rédigé automatiquement par un modèle d'IA.

% ================= 6. SECURITE =================
\chapter{Sécurité, qualité et bonnes pratiques}

Une part significative du stage a été consacrée à un travail d'audit, de sécurisation et de nettoyage du
code existant. Cet audit, mené par analyse statique du code (parcours de l'arbre syntaxique, compilation,
analyse des imports), a permis d'identifier et de corriger les points suivants~:

\begin{warnbox}
\textbf{Faille de sécurité critique n\textsuperscript{o}1 --- élévation de privilèges à l'inscription.} Le
formulaire d'inscription public acceptait initialement un champ \texttt{role} transmis par le client, ce qui
permettait à n'importe quel visiteur de créer un compte avec le rôle administrateur. Correction apportée~:
le rôle est désormais forcé côté serveur à « organisation » pour toute inscription publique~; la création
d'un compte administrateur nécessite l'exécution d'un script dédié, exécuté uniquement par un opérateur
autorisé.
\end{warnbox}

\begin{warnbox}
\textbf{Faille de sécurité critique n\textsuperscript{o}2 --- absence de vérification JWT sur les routes
d'administration.} Certains endpoints d'administration ne vérifiaient pas systématiquement la présence ou la
validité d'un jeton JWT. Un module centralisé (\texttt{dependencies.py}) a été créé et appliqué à l'ensemble
du routeur d'administration, garantissant qu'aucune route sensible n'est accessible sans authentification et
sans le rôle requis.
\end{warnbox}

Les autres actions de nettoyage et de fiabilisation menées durant le stage incluent~:

\begin{itemize}
  \item Suppression de code mort~: un service de recherche sémantique jamais référencé, un composant React
    jamais importé, et environ 470~Ko de jeux de données commis par erreur dans le dépôt~;
  \item Archivage d'une cinquantaine de scripts de diagnostic ponctuels et d'une trentaine de scripts de
    migration devenus obsolètes, déplacés vers des répertoires d'archives dédiés afin d'alléger la base de
    code active~;
  \item Correction d'un bug de configuration~: huit fichiers du frontend codaient en dur l'adresse
    \texttt{http://127.0.0.1:8000}, rendant tout déploiement hors environnement local impossible sans
    modification du code source. La résolution de l'URL de l'API est désormais entièrement pilotée par la
    variable d'environnement \texttt{VITE\_API\_URL}~;
  \item Correction d'un problème de dépendances~: les bibliothèques \texttt{axios} et
    \texttt{react-toastify}, utilisées par le frontend, étaient déclarées dans un \texttt{package.json} mal
    positionné à la racine du dépôt plutôt que dans celui du frontend, ce qui aurait provoqué un échec de
    build lors d'une installation propre des dépendances~;
  \item Suppression de dépendances inutilisées (\texttt{weasyprint}, \texttt{leaflet.heat}) et de 54 imports
    Python inutilisés dans le code backend~;
  \item Ajout de la conteneurisation complète du projet (voir chapitre~7).
\end{itemize}

Ce travail d'audit a permis de consolider la fiabilité et la sécurité de la plateforme avant son
déploiement, et constitue l'un des apports concrets les plus importants de ce stage.

% ================= 7. DEPLOIEMENT =================
\chapter{Déploiement et conteneurisation}

Afin de faciliter le déploiement et la reproductibilité de l'environnement, l'ensemble de la plateforme a
été conteneurisé à l'aide de Docker et Docker Compose. L'orchestration comprend~:

\begin{itemize}
  \item un service \texttt{db} basé sur l'image \texttt{pgvector/pgvector:pg16}, fournissant PostgreSQL avec
    l'extension de recherche vectorielle~;
  \item un service \texttt{backend}, construit à partir d'une image Python 3.12 slim, exécuté sous un
    utilisateur non privilégié (\texttt{appuser}) pour limiter la surface d'attaque en cas de compromission
    du conteneur~;
  \item un service \texttt{frontend}, construit en plusieurs étapes (build Vite sous Node 22, puis service
    des fichiers statiques via Nginx)~;
  \item deux services optionnels activables via des profils Docker Compose~: \texttt{elasticsearch} (profil
    \emph{search}) pour la recherche plein texte, et \texttt{ollama} (profil \emph{ai}) pour l'exécution
    locale du modèle de langage utilisé par le chatbot.
\end{itemize}

Cette approche par profils permet de démarrer une configuration minimale pour le développement courant
(\texttt{docker compose up}), tout en conservant la possibilité d'activer les fonctionnalités avancées à la
demande (\texttt{docker compose --profile search --profile ai up}).

% ================= 8. DIFFICULTES =================
\chapter{Difficultés rencontrées et solutions apportées}

\begin{longtable}{@{}p{6.5cm}p{7.5cm}@{}}
\toprule
\textbf{Difficulté} & \textbf{Solution apportée} \\
\midrule
\endhead
Base de données historiquement dénormalisée, mêlant données de test et données réelles &
  Refonte complète vers un schéma normalisé de 18 tables, documentée dans un schéma de base de données et un
  diagramme entité-association dédiés \\
Faille d'élévation de privilèges dans le flux d'inscription &
  Audit de sécurité ciblé, correction du rôle forcé côté serveur et mise en place d'une vérification JWT
  systématique \\
Adresse de l'API backend codée en dur, empêchant tout déploiement hors localhost &
  Externalisation de la configuration via la variable d'environnement \texttt{VITE\_API\_URL}, injectée au
  moment du build du frontend \\
Accumulation de scripts ponctuels et de code mort ralentissant la lisibilité du dépôt &
  Campagne de nettoyage et d'archivage, avec suppression des éléments non référencés et déplacement des
  scripts historiques vers des répertoires d'archives \\
Nécessité de supporter l'arabe en plus du français et de l'anglais, avec une direction d'écriture opposée &
  Mise en place d'une gestion dynamique de l'attribut \texttt{dir} (RTL/LTR) propagée jusqu'aux composants
  tiers (notifications, mise en page) \\
\bottomrule
\end{longtable}

% ================= 9. BILAN =================
\chapter{Bilan personnel et compétences acquises}

Ce stage m'a permis de développer et de consolider un ensemble de compétences techniques et
méthodologiques~:

\begin{itemize}
  \item Développement d'API REST avec FastAPI
  \item Modélisation de données relationnelles (SQLAlchemy)
  \item Authentification JWT et contrôle d'accès basé sur les rôles
  \item Développement frontend avec React et Vite
  \item Internationalisation et accessibilité (i18n, RTL)
  \item Recherche sémantique et embeddings vectoriels
  \item Intégration de modèles de langage (LLM) via Ollama
  \item Conteneurisation avec Docker et Docker Compose
  \item Audit de sécurité et analyse statique de code
  \item Documentation d'architecture (ADR)
\end{itemize}

Au-delà des aspects techniques, ce stage a été l'occasion de travailler sur un projet réel, déjà
partiellement développé, nécessitant une phase de compréhension et d'audit avant toute nouvelle
contribution — une situation représentative du travail sur des bases de code existantes en entreprise. La
rédaction d'un rapport d'audit structuré et la priorisation des corrections (sécurité avant confort, code
actif avant code historique) m'ont également permis de développer un sens critique vis-à-vis de la qualité
logicielle.

% ================= 10. CONCLUSION =================
\chapter{Conclusion générale et perspectives}

Le stage effectué au sein de l'AICTO sur le projet SARAI a permis de contribuer à la mise en place d'une
plateforme complète de recensement des initiatives d'intelligence artificielle dans la région arabe,
alliant un backend robuste (FastAPI, PostgreSQL, pgvector), une interface utilisateur riche et multilingue
(React, i18next), et des fonctionnalités avancées telles que la recherche sémantique et le chatbot
conversationnel. Le travail d'audit et de sécurisation mené en parallèle a permis d'améliorer
significativement la fiabilité, la sécurité et la maintenabilité du projet avant son déploiement.

\textbf{Perspectives d'évolution identifiées~:}

\begin{itemize}
  \item Migration du système de migrations de schéma vers un outil dédié tel qu'Alembic~;
  \item Mise en place d'une suite de tests automatisés (pytest, TestClient) et d'une intégration continue
    (CI)~;
  \item Introduction d'un mécanisme de \emph{refresh token} pour renforcer la gestion des sessions JWT~;
  \item Restriction de la configuration CORS via des variables d'environnement spécifiques à chaque
    environnement de déploiement~;
  \item Extraction progressive des services de recherche et d'IA en composants indépendants, si la charge ou
    l'équipe venait à croître.
\end{itemize}

\end{document}
